\documentclass[12pt,a4paper]{ctexart}
\usepackage[margin=1in]{geometry}
\usepackage{amsmath,amssymb,bm,mathtools}
\usepackage{graphicx}
\usepackage{booktabs}
\usepackage{array}
\usepackage{caption}
\usepackage{subcaption}
\usepackage{float}
\usepackage{algorithm}
\usepackage{algpseudocode}
\usepackage{xcolor}
\usepackage{listings}
\usepackage{hyperref}

\hypersetup{colorlinks=true,linkcolor=blue,citecolor=blue,urlcolor=blue}

\lstdefinestyle{matlabstyle}{
  language=Matlab,
  basicstyle=\ttfamily\small,
  keywordstyle=\color{blue},
  commentstyle=\color{teal},
  stringstyle=\color{purple},
  numbers=left,
  numberstyle=\tiny\color{gray},
  stepnumber=1,
  numbersep=8pt,
  breaklines=true,
  frame=single,
  showstringspaces=false
}

\title{单比特量化约束下 MIMO-OFDM 通感一体化系统\\DFT 基高精度角度估计组会汇报}
\author{汇报人：\underline{请填写姓名} \quad 单位：\underline{请填写单位}}
\date{2026年4月}

\begin{document}
\maketitle

\begin{abstract}
本文围绕“单比特量化条件下如何实现高精度角度估计”展开。针对 MIMO-OFDM 通感一体化（ISAC）系统中高分辨 ADC 成本和功耗过高的问题，采用复单比特 ADC 对回波进行量化，分析其带来的幅度信息丢失、量化噪声抬升、通信符号污染谱峰等关键退化机制。在此基础上，构建“幅度补偿—空间维DFT分箱—自适应缩放因子消除通信符号—角度谱累积—CFAR与插值精化”的完整算法链路。结合 MATLAB 实现（\texttt{run\_demo\_1bit\_dft.m}、\texttt{joint\_arv\_estimator.m}、\texttt{run\_mc\_eval\_1bit\_angle.m}）给出可复现实验模板，并提供论文公式与代码步骤的一一映射。结果展示建议表明：Bussgang 补偿与谱峰精化可显著提升 1-bit 场景测角稳健性与精度，为低功耗终端部署提供工程可行路径。
\end{abstract}

\section{研究背景与问题定义}
ISAC 通过通信与感知频谱/硬件复用，成为 6G 车联网与智能终端的重要技术路线。MIMO-OFDM 具备成熟标准与良好多径鲁棒性，但大规模阵列与宽带配置使接收数据量巨大，导致高精度 ADC 在功耗与成本上难以落地。单比特 ADC 仅保留符号信息，可显著降低硬件复杂度，但会引入三类核心挑战：
\begin{itemize}
  \item \textbf{幅度信息丢失}：无法直接进行传统幅度域干扰消除；
  \item \textbf{量化噪声抬升}：角度谱主瓣变钝、旁瓣抬高；
  \item \textbf{通信符号污染}：随机符号在角域造成伪峰，影响峰值定位。
\end{itemize}
因此，本工作目标是：在 1-bit 量化下实现低复杂度、可解释、可复现的高精度角度估计。

\section{系统模型与物理机理}
对第 $m$ 根接收天线、第 $i$ 个子载波、第 $l$ 个 OFDM 符号，回波建模为
\begin{equation}
  y(m,i,l)=\sum_{q=1}^{Q}\beta_q\,\mathbf{a}^H(\theta_q)\mathbf{x}_i[l]\,e^{j m\omega_a(\theta_q)}e^{j i\omega_r(R_q)}e^{j l\omega_v(v_q)}+z(m,i,l).
\end{equation}
其中角度数字频率 $\omega_a(\theta_q)=-\frac{2\pi d_r}{\lambda_c}\sin\theta_q$，距离频率 $\omega_r(R_q)=-\frac{4\pi\Delta f}{c}R_q$，速度频率 $\omega_v(v_q)=\frac{4\pi T f_c}{c}v_q$。关键结论是：角度、距离、速度分别耦合在\textbf{空间维}、\textbf{快时间维}、\textbf{慢时间维}，因此角度估计可优先在空间维执行 DFT。

\section{单比特量化建模与退化分析}
复单比特量化表达为
\begin{equation}
  y_b=\mathcal{Q}(y)=\operatorname{sign}(\Re\{y\})+j\,\operatorname{sign}(\Im\{y\}).
\end{equation}
在 Bussgang decomposition 下近似写成
\begin{equation}
  y_b = k y + n_q, \quad k=\frac{2}{\sqrt{\pi}},
\end{equation}
其中 $n_q$ 为与原信号统计不相关的量化噪声。该模型解释了 1-bit 条件下谱峰钝化和噪底上升。

\section{方法：单比特空间维 DFT 高精度角度估计}
\subsection{步骤1：幅度补偿}
利用 $k$ 对量化信号进行比例恢复，减轻量化收缩效应：
\begin{equation}
  \tilde{y}_b = k\,y_b.
\end{equation}

\subsection{步骤2：空间维 DFT 分箱}
对空间维执行 $N_a$ 点 DFT：
\begin{equation}
  Y_{b,i,l}(n_a)=\frac{1}{M_{rx}}\sum_{m=0}^{M_{rx}-1}\tilde{y}_b(m,i,l)e^{-j m\tilde\omega_a(n_a)}.
\end{equation}
由于距离/速度项与 $m$ 无关，它们作为复常数不改变角域分箱机理。

\subsection{步骤3：自适应缩放因子消除通信符号}
定义
\begin{equation}
  y_{b,i,l}(n_a)=\frac{Y_{b,i,l}(n_a)}{\alpha_{b,n_a}\,\mathbf{a}^H(\theta_{n_a})\mathbf{x}_i[l]}.
\end{equation}
通过功率守恒约束求解 $\alpha_{b,n_a}$，避免直接相除导致角度箱功率畸变。

\subsection{步骤4：角度谱构造与峰值精化}
沿 $i,l$ 维累积得角度功率谱
\begin{equation}
  Y_b(n_a)=\frac{1}{N_sL}\sum_{i=0}^{N_s-1}\sum_{l=0}^{L-1}|y_{b,i,l}(n_a)|^2.
\end{equation}
随后采用 CA-CFAR（Constant False Alarm Rate）抑制虚警，并用抛物线插值（parabolic interpolation）提升亚栅格定位精度。

\subsection{步骤5：角度反演}
\begin{equation}
  \hat\theta = \arcsin\!\left(-\frac{\hat n_a\lambda_c}{d_r N_a}\right).
\end{equation}

\begin{algorithm}[H]
\caption{单比特空间维 DFT 角度估计流程}
\begin{algorithmic}[1]
\Require $y_b(m,i,l)$，系统参数 $M_{rx},N_a,d_r,\lambda_c$，已知发射符号 $x$
\Ensure $\hat\theta$
\State 幅度补偿：$\tilde y_b \leftarrow k y_b$
\State 空间维 DFT：计算 $Y_{b,i,l}(n_a)$
\State 求解 $\alpha_{b,n_a}$ 并消除 $\mathbf{a}^H(\theta_{n_a})\mathbf{x}_i[l]$
\State 跨 $i,l$ 维累积得到 $Y_b(n_a)$
\State CA-CFAR + 抛物线插值，得到 $\hat n_a$
\State 角度反演输出 $\hat\theta$
\end{algorithmic}
\end{algorithm}

\section{论文到代码映射（重点汇报页）}
\subsection*{1) 参数与物理含义：\texttt{run\_demo\_1bit\_dft.m}}
建议展示：\texttt{mu\_3gpp, Delta\_f, fc, lambda\_c, dt, dr, Na/Nr/Nv, T} 的定义。强调“参数-物理意义-公式变量”的对应关系。

\subsection*{2) 主算法实现：\texttt{joint\_arv\_estimator.m}}
建议按代码注释中的 Step 1--Step 5 逐页对应：
\begin{itemize}
  \item Step 1：\texttt{use\_bussgang} 开关与幅度补偿；
  \item Step 2：\texttt{fft(...,Na,1)} 实现空间维 DFT；
  \item Step 3：逐角度箱构造 \texttt{denom} 与 \texttt{alpha}，完成通信符号消除；
  \item Step 4：\texttt{fft2} 构造距离-速度域；
  \item Step 5：\texttt{asin}、\texttt{R\_hat}、\texttt{v\_hat} 完成参数反演。
\end{itemize}

\subsection*{3) 蒙特卡洛评测：\texttt{run\_mc\_eval\_1bit\_angle.m}}
建议展示三组对比：\texttt{full}、\texttt{1bit\_no\_bg}、\texttt{1bit\_bg}，并突出 RMSE、中位绝对误差、离群率三指标。

\section{实验设计与结果展示模板}
\subsection{可复现实验设置}
单目标场景：$\theta=10^\circ$，$R=60\,\mathrm{m}$，$v=20\,\mathrm{m/s}$，SNR=15 dB；
默认规模可选 balanced：$M_{rx}=16,N_{tx}=16,N_s=256,L=64,N_a=3M_{rx}$。

\subsection{建议图表}
\begin{figure}[H]
  \centering
  \includegraphics[width=0.8\linewidth]{figures/angle_spectrum_placeholder.pdf}
  \caption{角度谱示意图（请替换为程序输出图）}
  \label{fig:angle_spectrum}
\end{figure}

\begin{figure}[H]
  \centering
  \includegraphics[width=0.8\linewidth]{figures/rmse_snr_placeholder.pdf}
  \caption{角度 RMSE 随 SNR 变化（full / 1bit-noBG / 1bit-BG）}
  \label{fig:rmse_snr}
\end{figure}

\begin{table}[H]
\centering
\caption{不同算法配置性能对比（示例模板）}
\begin{tabular}{c c c c}
\toprule
SNR(dB) & RMSE(full) & RMSE(1bit-noBG) & RMSE(1bit-BG) \\
\midrule
-5 & -- & -- & -- \\
0  & -- & -- & -- \\
5  & -- & -- & -- \\
10 & -- & -- & -- \\
15 & -- & -- & -- \\
\bottomrule
\end{tabular}
\end{table}

\section{组会讲解提纲（5--8分钟）}
\begin{enumerate}
  \item 研究动机（1分钟）：为什么要 1-bit，代价是什么；
  \item 模型与机理（1分钟）：角-距-速维度解耦；
  \item 方法链路（2分钟）：补偿、分箱、消符号、谱构造、精化；
  \item 代码映射（1分钟）：论文公式如何落地到函数与开关；
  \item 结果与结论（1分钟）：BG 与精化带来的收益；
  \item 讨论（1--2分钟）：多目标、近场、阈值优化等后续方向。
\end{enumerate}

\section{局限性与改进方向}
\begin{itemize}
  \item 当前以单目标/理想阵列为主，需扩展到多目标近间隔场景；
  \item 远场窄带近似下效果较好，需评估近场与宽带波束偏斜影响；
  \item 量化阈值固定，后续可研究自适应阈值与非对称量化；
  \item 可补充与 MUSIC/ESPRIT 在 1-bit 条件下的统一对比；
  \item 推进硬件在环（HIL）验证，评估实际链路鲁棒性。
\end{itemize}

\section{结论}
本文给出一条从物理建模到工程实现的 1-bit 高精度测角路径：通过空间维解耦把角度估计问题转化为低复杂度频域峰值检测，并以 Bussgang 补偿与自适应缩放因子缓解量化与通信符号干扰。结合 CFAR 与插值后，在低功耗前提下可获得可用的角度估计精度，具备面向终端部署的实践价值。

\section*{附：代码展示页（可直接嵌入）}
\noindent 若希望在 PDF 中直接展示源码片段，可取消注释并调整行号：

% \lstset{style=matlabstyle}
% \lstinputlisting[firstline=1,lastline=120,caption={run\_demo\_1bit\_dft.m 参数与仿真设置}]{run_demo_1bit_dft.m}
% \lstinputlisting[firstline=1,lastline=220,caption={joint\_arv\_estimator.m 主算法实现}]{joint_arv_estimator.m}
% \lstinputlisting[firstline=1,lastline=220,caption={run\_mc\_eval\_1bit\_angle.m 蒙特卡洛评测}]{run_mc_eval_1bit_angle.m}

\begin{thebibliography}{99}
\bibitem{isac_overview}
F. Liu, Y.-F. Liu, A. Li, C. Masouros, Y. C. Eldar, and S. Cui,
``Integrated sensing and communications: Toward dual-functional wireless networks,''
\emph{IEEE Journal on Selected Areas in Communications}, 2022.

\bibitem{1bit_adc}
J. Mo and R. W. Heath,
``High SNR capacity of millimeter wave MIMO systems with one-bit quantization,''
\emph{Information Theory and Applications Workshop}, 2014.

\bibitem{bussgang}
J. J. Bussgang,
``Crosscorrelation functions of amplitude-distorted Gaussian signals,''
\emph{MIT Research Laboratory of Electronics, Technical Report}, 1952.

\bibitem{cfar}
H. Rohling,
``Radar CFAR thresholding in clutter and multiple target situations,''
\emph{IEEE Transactions on Aerospace and Electronic Systems}, 1983.
\end{thebibliography}

\end{document}
